\documentclass[12pt,oneside]{book}

% ==========================================================
% METADATOS
% ==========================================================
\newcommand{\universidad}{Universidad de Buenos Aires (UBA)}
\newcommand{\facultad}{Facultad de Derecho}
\newcommand{\materia}{Práctica Profesional}
\newcommand{\titulo}{Recursos y apelación en el proceso civil}
\newcommand{\autor}{Isaac Edgar Camacho Ocampo}
\newcommand{\anio}{2026}

% ==========================================================
% PAQUETES
% ==========================================================
\usepackage[utf8]{inputenc}
\usepackage[T1]{fontenc}
% Se evita babel en español porque no está instalado en este entorno;
% esto permite compilar el documento y generar el índice.
\renewcommand{\contentsname}{Índice}

\usepackage[
    top=2.5cm,
    bottom=2.5cm,
    left=3cm,
    right=2.5cm,
    headheight=15pt
]{geometry}

\usepackage{amsmath}
\usepackage{amssymb}
\usepackage{mathtools}

\usepackage{graphicx}
\usepackage{float}
\usepackage{caption}
\usepackage{subcaption}

\usepackage{booktabs}
\usepackage{array}
\usepackage{longtable}
\usepackage{tabularx}

\usepackage{enumitem}

\usepackage{xcolor}
\usepackage[most]{tcolorbox}

\usepackage{fancyhdr}
\usepackage{ragged2e}

\usepackage[
    colorlinks=true,
    linkcolor=blue!50!black,
    citecolor=green!40!black,
    urlcolor=blue!60!black,
    pdftitle={\titulo},
    pdfauthor={\autor},
    pdfsubject={Apuntes universitarios},
    bookmarks=true
]{hyperref}

% ==========================================================
% CONFIGURACIÓN
% ==========================================================
\graphicspath{
    {./img/}
    {./graficos/}
    {./figuras/}
}

\setcounter{tocdepth}{2}

\setlength{\parindent}{0pt}
\setlength{\parskip}{0.5em}

\setlist[itemize]{
    topsep=0.3em,
    itemsep=0.2em
}

\setlist[enumerate]{
    topsep=0.3em,
    itemsep=0.2em
}

\clubpenalty=10000
\widowpenalty=10000

% ==========================================================
% COMANDOS
% ==========================================================
\newcommand{\abs}[1]{\left|#1\right|}
\newcommand{\norm}[1]{\left\lVert#1\right\rVert}
\newcommand{\vect}[1]{\mathbf{#1}}
\newcommand{\deriv}[2]{\frac{d#1}{d#2}}
\newcommand{\pderiv}[2]{\frac{\partial#1}{\partial#2}}

% ==========================================================
% ENTORNOS / CAJAS
% ==========================================================
\newtcolorbox{definition}[1]{
    enhanced,
    breakable,
    title={Definición: #1},
    fonttitle=\bfseries,
    colback=blue!3,
    colframe=blue!50!black,
    boxrule=0.8pt,
    arc=2mm
}

\newtcolorbox{important}{
    enhanced,
    breakable,
    title={Importante},
    fonttitle=\bfseries,
    colback=yellow!8,
    colframe=orange!70!black,
    boxrule=0.8pt,
    arc=2mm
}

\newtcolorbox{example}[1]{
    enhanced,
    breakable,
    title={Ejemplo: #1},
    fonttitle=\bfseries,
    colback=green!4,
    colframe=green!40!black,
    boxrule=0.8pt,
    arc=2mm
}

\newtcolorbox{formula}[1]{
    enhanced,
    breakable,
    title={Fórmula / Regla: #1},
    fonttitle=\bfseries,
    colback=purple!4,
    colframe=purple!50!black,
    boxrule=0.8pt,
    arc=2mm
}

\newtcolorbox{exercise}[1]{
    enhanced,
    breakable,
    title={Ejercicio: #1},
    fonttitle=\bfseries,
    colback=gray!5,
    colframe=gray!60!black,
    boxrule=0.8pt,
    arc=2mm
}

\newtcolorbox{note}{
    enhanced,
    breakable,
    title={Nota},
    fonttitle=\bfseries,
    colback=white,
    colframe=gray!50,
    boxrule=0.6pt,
    arc=2mm
}

% ==========================================================
% CONFIGURACIÓN FINAL
% ==========================================================
\pagestyle{fancy}
\fancyhf{}

\fancyhead[L]{\nouppercase{\leftmark}}
\fancyhead[R]{\materia}

\fancyfoot[C]{\thepage}

\renewcommand{\headrulewidth}{0.4pt}
\renewcommand{\footrulewidth}{0pt}

\fancypagestyle{plain}{
    \fancyhf{}
    \fancyfoot[C]{\thepage}
    \renewcommand{\headrulewidth}{0pt}
}

\renewcommand{\arraystretch}{1.25}

% ==========================================================
\begin{document}

% ==========================================================
% PORTADA
% ==========================================================
\begin{titlepage}

\thispagestyle{empty}

\begin{center}

\vspace*{1cm}

{\large \textsc{\universidad}}

\vspace{0.2cm}

{\large \textsc{\facultad}}

\vspace{3cm}

{\Huge \textbf{\materia}}

\vspace{1cm}

{\LARGE \titulo}

\vspace{0.8cm}

{\small \textit{Estudiar con constancia es el camino para convertir lo aprendido en criterio propio.}}

\vspace{1.2cm}

\rule{0.70\textwidth}{0.5pt}

\vspace{0.5cm}

{\large Apuntes de estudio}

\vspace{0.5cm}

\rule{0.70\textwidth}{0.5pt}

\vfill

{\large \autor}

\vspace{0.3cm}

{\large \anio}

\end{center}

\vfill

\begin{flushleft}
\textit{Este es un modesto aporte a los estudiantes de ``Práctica Profesional'' no pretende sustituir el material oficial de la materia.}
\end{flushleft}

\end{titlepage}

% ==========================================================
% ÍNDICE
% ==========================================================
\frontmatter
\tableofcontents
\mainmatter

% ==========================================================
% CONTENIDO
% ==========================================================

% ----------------------------------------------------------
\chapter[Visión general]{Visión general de los temas}
% ----------------------------------------------------------

\section{Temas tratados}

Los contenidos combinan una reflexión sobre el \textbf{recurso extraordinario} y la doctrina de la arbitrariedad con un taller práctico sobre la construcción de una \textbf{expresión de agravios}. Dicha expresión de agravios debe contener una crítica concreta y razonada (art.~265 del Código Procesal). El material incluye el análisis de un caso práctico sobre ruidos de una obra lindera (art.~1973 del Código Civil y Comercial), los límites de la \textbf{prueba en segunda instancia} (art.~260) y la \textbf{apelación de honorarios} (art.~244 del Código Procesal y Ley de Aranceles 27.423): plazos, fundamentación, mora e intereses.

\section{Relación entre los temas}

Los temas forman un único recorrido: \textbf{reclamar, revisar, fundar y cobrar}. Puede ilustrarse con el caso de un vecino que realiza una obra con ruidos de la mañana a la noche:

\begin{enumerate}
    \item Quien reclama tiene derechos y la Constitución no es un adorno: se vive reclamando.
    \item El primer juez rechaza el pedido porque «no hay mediciones». Para pedir su revisión es necesario explicar \emph{qué dijo el juez, por qué está equivocado y qué norma lo muestra}: esa es la crítica concreta y razonada.
    \item En la segunda revisión no se arma un juicio nuevo: se discute lo ya actuado, y un escrito sin firma «no existe».
    \item El trabajo del abogado tiene un precio: los honorarios también se pueden apelar, con plazos y riesgos de mora.
\end{enumerate}

\section{Utilidad para la vida profesional}

\begin{note}
Estos contenidos resultan directamente útiles por cuatro razones concretas:
\begin{enumerate}
    \item \textbf{Evita perder un caso que se tenía ganado.} Un recurso sin crítica concreta y razonada se declara desierto: el tribunal no suple lo que la parte omite.
    \item \textbf{Protege el ingreso.} Los honorarios tienen plazos de 5 y 10 días, y la mora genera intereses retroactivos; saber apelarlos y cobrarlos es vivir de la profesión.
    \item \textbf{Enseña a decidir con estrategia.} Fundar o no fundar, cuándo hacerlo y qué riesgos abre cada opción son decisiones que impactan en el cliente.
    \item \textbf{Se aplica fuera de los tribunales.} Argumentar con precisión («qué dijo, por qué está mal, qué norma lo respalda»), convivir con la inteligencia artificial como herramienta y reclamar sin dar nada por supuesto sirve para un reclamo al consorcio, a un empleador o a una empresa.
\end{enumerate}
\end{note}

% ----------------------------------------------------------
\chapter[Recurso extraordinario y Constitución]{Del recurso extraordinario a la Constitución vivida}
% ----------------------------------------------------------

\section{Arbitrariedad y recurso extraordinario}

\subsection{Qué del recurso extraordinario se ve en las apelaciones comunes}

Un elemento propio del recurso extraordinario que reaparece en los recursos comunes es la \textbf{doctrina de la arbitrariedad}, muy discutida y con gran presencia en la práctica. Se trata de una línea de fallos de la Corte: quien busque jurisprudencia encontrará más fallos de arbitrariedad que de derecho federal escrito. Esto se explica porque la materia federal no es tan vasta y porque, en las materias de derecho común, la Corte no declara permanentemente que estén excluidas del recurso extraordinario. En determinados casos, la Corte admite que se conozcan por esta vía.

\begin{definition}{Doctrina de la arbitrariedad}
Línea jurisprudencial de la Corte que permite abrir la vía extraordinaria aun cuando el caso sea de derecho común.
\end{definition}

\subsection{La paradoja de pedir arbitrariedad al mismo tribunal}

El recurso extraordinario debe interponerse ante el mismo tribunal que dictó la sentencia que se reputa arbitraria. Ello resulta en ocasiones contradictorio: es difícil que un tribunal reconozca que dictó una sentencia arbitraria. Puede ocurrir que el tribunal, aun negando la arbitrariedad, considere que no le corresponde resolverla y que debe hacerlo el Tribunal Superior. Aun así, el planteo está así previsto y a veces no queda otro remedio.

\subsection{Cuándo se abre la cuestión federal en un caso de derecho común}

En un reclamo de derecho común, la pregunta es dónde está la norma federal. En principio no la hay. Sin embargo, si el juez, para resolver la cuestión de derecho común, declaró la inconstitucionalidad de una norma, es posible que se abra la cuestión federal.

\begin{important}
Hay más fallos de arbitrariedad que de derecho federal escrito porque la materia federal es acotada. Una vía concreta para abrir la cuestión federal es que el juez haya declarado la inconstitucionalidad de una norma.
\end{important}

\section{Ley 48, artículo 14 y cuestión federal}

Para conocer las hipótesis del recurso extraordinario hay que leer la \textbf{ley 48} y los supuestos de su \textbf{artículo 14}. Esa ley fue sancionada en los albores de la organización nacional y es allí donde se contemplaron las hipótesis del recurso extraordinario.

% TODO: contenido incompleto o ambiguo en las notas originales (el texto de la ley 48, art. 14, no fue leído en clase)
\begin{note}
La norma no fue leída en clase: se remite a su lectura como fuente de las hipótesis del recurso extraordinario federal.
\end{note}

\subsection{La evolución de la «cuestión federal» en la práctica profesional}

Cuando se comenzó a ejercer la profesión, prácticamente no se mencionaba la cuestión federal: si se hablaba de una cuestión de derecho común, no se veía dónde estaba la cuestión federal. Con el tiempo esto ha ido creciendo y ganando envergadura, hasta el punto de que los derechos constitucionales se ejercen desde el primer momento en que se acude a un juzgado.

\section{Constitución viva, derechos humanos y participación popular}

\subsection{«Vivir la Constitución»}

Estar en una situación infraconstitucional no significa que la Constitución sea irrelevante. La Constitución es algo para ser \textbf{vivido}: no es una entelequia ubicada arriba, sino que debe practicarse y aplicarse. Por eso se menciona siempre alguna norma de la Constitución, desde el acceso a la justicia, y no conviene quedarse en una separación entre lo constitucional y lo infraconstitucional.

\subsection{Derecho internacional: monismo y dualismo}

Cada vez más se aplican tratados de derechos humanos («lo convencional»). Durante mucho tiempo se vivió como si el orden interno y el orden internacional fueran dos cosas distintas, prácticamente separadas, según el debate entre las doctrinas del \textbf{monismo} y del \textbf{dualismo}. Hoy se aplican muchas de esas normas internacionales.

\subsection{Normas operativas o programáticas: el juicio por jurados}

Respecto de la discusión constitucional sobre si las normas son operativas o programáticas, se sostiene que son \textbf{operativas}, aunque existan normas que han tardado en implementarse. El ejemplo es el juicio por jurados: pasaron muchos años hasta que se lo tuvo, y aún no existe en todas las ramas, como prevé el artículo 24 de la Constitución, por lo que debería haber juicio por jurados también en las temáticas civiles. Esto significa, de algún modo, que la gente esté más cerca de la aplicación del derecho y de lo que se llama justicia.

Como comparación, en Estados Unidos hay modalidades de juicio por jurados no solamente en cuestiones criminales, por ejemplo en las acciones de clase. En el país no existen, pero se está trabajando en esa dirección.

% TODO: contenido incompleto o ambiguo en las notas originales (el texto del art. 24 CN no fue leído en clase)
\begin{note}
El artículo 24 de la Constitución Nacional se cita para señalar que la Constitución prevé el juicio por jurados y que aún no está implementado «en todas las ramas», incluidas las materias civiles. Su texto no fue leído.
\end{note}

\subsection{Elección de los jueces y participación popular}

Sobre la elección de los jueces se plantea quién elige al Poder Judicial: no lo elige la gente, sino que lo hace un mecanismo indirecto. Antes de que existiera el Consejo de la Magistratura todo era más discrecional; en la opinión expresada, el mecanismo sigue teniendo un componente discrecional, porque es una elección a través de representantes que la gente no vota. Los abogados tienen el privilegio de votar representantes (recientemente se votó para el Consejo de la Magistratura de la Nación) y los jueces también votan. La gente, en cambio, vota solo a través de los políticos que participan, ya que hay representantes del Poder Ejecutivo y del Poder Legislativo, pero todo de manera indirecta. La participación popular en los mecanismos de juzgamiento permite que la gente esté más cerca y entienda un poco más.

\section{Acceso a la justicia e inteligencia artificial}

\subsection{El expediente es inaccesible para la gente}

Hoy una persona tiene la posibilidad de ver su expediente, pero no sabe cómo hacerlo. Si no es a través de los abogados, nadie puede acceder a un expediente. Antes tampoco era maravilloso: la gente iba al juzgado y se le mostraba el expediente, pero muchas veces no lo entendía. Al preguntar al empleado de mesa de entradas, este respondía que no puede asesorar y remitía a un abogado o al patrocinio jurídico gratuito de la facultad.

A los consultantes no se les muestra el expediente porque no lo entienden; lo que puede tenerse es una carpeta impresa. A la gente le cuesta muchísimo leer, y si cuesta a los abogados y a los adultos mayores, el expediente resulta inaccesible. Por un lado la tecnología dinamiza, pero se trata de una de las externalidades no queridas del desarrollo, que puede ser un retraso o un perjuicio porque la gente no puede acceder.

\subsection{La inteligencia artificial y la profesión}

Hay consultantes que llegan diciendo que subieron el expediente a la inteligencia artificial y que esta les indica lo que hay que hacer; la profesión tiene que convivir con eso. Se relata el caso de un contador que consultó a una herramienta de IA cómo enviar una carta documento a la encargada de un consorcio que debía jubilarse; la herramienta le indicó prácticamente todo el recorrido. Sin embargo, hay personas criteriosas que igualmente consultan al abogado, porque si algo sale mal existe responsabilidad profesional y los consorcistas preguntarán si se consultó con un abogado. Por algo es una \textbf{herramienta}.

Además, la herramienta puede aportar ideas que no se habían ocurrido de inmediato, y también el abogado puede usarla para pensar el recorrido. Se mencionaron otros aspectos:

\begin{itemize}
    \item Es posible «entrenarla» y trabajar en conjunto, aunque uno puede volverse dependiente.
    \item Es muy rápida, pero tiene costo: hay que pagar, y hay que «alimentarla».
    \item Respecto de la seguridad, las personas que trabajan en la aplicación pueden verlo todo, por lo que se advirtió que no hay que dar datos personales.
    \item Antes eran herramientas exclusivas y caras; ahora accede cualquiera por un abono. Se la usa, por ejemplo, para cálculos de liquidaciones o en temas de pericia contable.
    \item Los profesionales enfrentan permanentemente nuevos desafíos, como en su momento los cursos de informatización del sistema.
\end{itemize}

\begin{important}
La inteligencia artificial es una herramienta: acelera el trabajo y puede aportar un recorrido, pero la responsabilidad profesional, el criterio y la reflexión profunda siguen siendo del abogado. Tiene costo, exige cuidado con los datos personales y plantea el desafío permanente de la actualización.
\end{important}

\section{Cuadro Cornell del capítulo}

\begin{longtable}{>{\raggedright\arraybackslash}p{0.27\textwidth}>{\raggedright\arraybackslash}p{0.67\textwidth}}
\toprule
\textbf{Preguntas / palabras clave} & \textbf{Notas / respuestas} \\
\midrule
\endfirsthead
\toprule
\textbf{Preguntas / palabras clave} & \textbf{Notas / respuestas} \\
\midrule
\endhead

\textbf{¿Cuándo se abre el recurso extraordinario en un caso de derecho común?} &
Por la doctrina de la arbitrariedad, una línea de fallos de la Corte más abundante que la de derecho federal escrito, porque la materia federal es acotada. También cuando el juez declaró la inconstitucionalidad de una norma, lo que abre la cuestión federal (ley 48, art.~14). \\[0.6em]

\textbf{¿Qué paradoja presenta pedir la arbitrariedad?} &
El recurso se interpone ante el mismo tribunal que dictó la sentencia que se reputa arbitraria, que difícilmente reconocerá su arbitrariedad; puede además considerar que debe resolverlo el Tribunal Superior. \\[0.6em]

\textbf{¿Qué significa «vivir la Constitución»?} &
Aplicarla y practicarla desde el primer momento en que se acude a un juzgado, y no solo citarla: acceso a la justicia, tratados de derechos humanos, normas operativas (por ejemplo, juicio por jurados, art.~24 CN) y participación popular. \\[0.6em]

\textbf{¿Qué debate plantean el monismo y el dualismo?} &
Si el orden interno y el internacional son dos cosas separadas o no. Durante mucho tiempo se los vivió como separados; hoy se aplican muchas normas internacionales. \\[0.6em]

\textbf{¿Qué rol tiene la inteligencia artificial?} &
Es una herramienta: acelera y sugiere recorridos, pero tiene costo, riesgos respecto de los datos personales y no reemplaza la responsabilidad ni el criterio profesional del abogado. \\

\midrule
\multicolumn{2}{>{\raggedright\arraybackslash}p{0.94\textwidth}}{%
\textbf{Resumen:} la doctrina de la arbitrariedad permite abrir la vía extraordinaria en casos de derecho común, y la declaración de inconstitucionalidad abre la cuestión federal. La Constitución debe vivirse y aplicarse en cada escrito. El expediente es difícilmente accesible para la gente y la inteligencia artificial es una herramienta útil, pero la responsabilidad y el criterio siguen siendo del abogado.
} \\
\bottomrule
\end{longtable}

% ----------------------------------------------------------
\chapter[La expresión de agravios]{La expresión de agravios: crítica concreta y razonada}
% ----------------------------------------------------------

\section{Crítica concreta y razonada}

\subsection{Del principio constitucional a la lesión concreta}

Es importante vivir las normas constitucionales, aunque cada uno las aplique de acuerdo con su mirada. Una cosa es decir «se violó el derecho de defensa en juicio» (en rigor, la garantía de la defensa en juicio) y otra es decir concretamente: «por haber decidido de tal manera, se lesiona este precepto constitucional». Esto es desafiante, porque hay que saber explicarlo y reflejarlo en el escrito y en el recurso.

\subsection{El artículo 265}

El artículo 265 del Código Procesal Civil y Comercial de la Nación «manda» que la expresión de agravios contenga una \textbf{crítica concreta y razonada} de lo resuelto.

\begin{definition}{Crítica concreta y razonada (art.~265 CPCCN)}
Exigencia de que la expresión de agravios señale concretamente qué parte de lo resuelto es equivocada y por qué, mostrando cómo y por qué una decisión lesiona un precepto determinado.
\end{definition}

% TODO: contenido incompleto o ambiguo en las notas originales (el texto del art. 265 no fue leído en clase; la fórmula es la que utiliza el docente)
\begin{note}
El texto del artículo 265 no fue leído en clase; la fórmula «crítica concreta y razonada» es la que utiliza el docente.
\end{note}

\section{El objeto del escrito}

En el caso práctico, el capítulo «Objeto» de la expresión de agravios fue redactado de la siguiente manera.

\begin{example}{Capítulo «Objeto» (caso práctico)}
\textbf{I. OBJETO.} Se expresan agravios contra la sentencia definitiva dictada el 28 de octubre de 2026 por el Juzgado Nacional de Primera Instancia en lo Civil n.º 33, que rechazó íntegramente la demanda e impuso las costas a la parte actora. Se solicita su revocación parcial respecto del codemandado Alfredo Bustos, en cuanto desestimó la pretensión resarcitoria fundada en las molestias provenientes de la obra lindera y el daño extrapatrimonial atribuido a ellas. Se pide asimismo la revisión de las costas con arreglo al resultado del recurso o, subsidiariamente, la aplicación de la excepción prevista en el artículo 68, segundo párrafo, del Código Procesal Civil y Comercial de la Nación. La crítica se formula en los términos del artículo 265 del mismo Código.
\end{example}

El artículo 68, segundo párrafo, del CPCCN fue citado en el objeto como la excepción al principio de imposición de costas al vencido (aplicación subsidiaria); su texto no fue leído en clase.

\subsection{Observaciones sobre el objeto}

A veces, en los objetos se adelantan peticiones que formarían parte del petitorio. No está mal, pero son cuestiones de forma. El objeto puede limitarse a expresar «vengo a expresar agravios», y extenderse sobre lo que se pretende forma parte de otro capítulo. No se trata de ser purista en la redacción: lo que se busca es transmitir en ese capítulo \textbf{el propósito del escrito}, que se desarrollará en los apartados siguientes, es decir, la crítica concreta y razonada.

\section{Estructura de los agravios}

Es habitual estructurar la expresión de agravios con los encabezados «primer agravio», «segundo agravio», o directamente con el nombre de los agravios. En el caso práctico se procedió así:

\begin{itemize}
    \item \textbf{Primer agravio:} valoración insuficiente de los ruidos y de la prueba testimonial.
    \item \textbf{Segundo agravio:} examen insuficiente del daño extrapatrimonial.
    \item No hay tercer agravio; luego se tratan costas y honorarios.
\end{itemize}

Se consideró que esta forma es una posibilidad válida, y quizás la que más conviene, aunque puede haber otras.

\section{Autosuficiencia del recurso: parafrasear la sentencia}

Muchos memoriales o expresiones de agravios comienzan analizando lo que dice la sentencia: qué dijo el juzgado, cómo se resolvió el conflicto. El motivo es que el recurso tiene que tener \textbf{autosuficiencia}: su lectura debe permitir entender todo el complejo de la causa, porque en la sentencia se supone que se condensó todo lo ocurrido en el proceso, y esa es la plataforma del recurso: qué decidió el juez y cómo lo decidió.

No se trata de copiar y pegar la sentencia en el recurso, sino de \textbf{parafrasearla}, con una referencia lo más textual posible, y recién después criticarla. Los considerandos son las partes del razonamiento de la sentencia, de modo que el planteo sigue este esquema: «Considerando tal, se dijo tal cosa. Eso es equivocado. Es equivocado por tales razones». Eso es hacer una crítica concreta. Aunque la cámara vea la sentencia, esta debe mencionarse concretamente en el recurso. Esto es lo que reclama la normativa con la expresión «crítica concreta y razonada».

\begin{important}
Una buena expresión de agravios es autosuficiente: primero reconstruye qué decidió el juez y cómo lo decidió (paráfrasis lo más textual posible, identificando el considerando) y recién después lo critica, señalando por qué es equivocado. Mencionar concretamente la sentencia es lo que convierte la queja en una crítica concreta y razonada.
\end{important}

\section{Primer agravio: valoración de la prueba}

El primer agravio del caso práctico dice lo que dijo el juez en la sentencia y luego lo critica. Se reproduce su desarrollo.

\begin{example}{Primer agravio: valoración insuficiente de los ruidos y de la prueba testimonial}
\textbf{La sentencia.} El considerando 9 reconoce que los testimonios aportan referencias coincidentes sobre la percepción de ruidos provenientes de la obra y su carácter molesto. Sin embargo, destaca que no se incorporaron mediciones objetivas que permitieran determinar su intensidad y, en consecuencia, concluye que no se encontraba acreditado que pudieran exceder la normal tolerancia.

\textbf{El agravio.} No consiste en sostener que el fallo haya establecido expresamente que las mediciones acústicas constituyen el único medio de prueba admisible. Consiste en que, pese a reconocer la existencia de testimonios coincidentes y relevantes, la sentencia no desarrolla suficientemente por qué las circunstancias concretas que surgen de esas declaraciones resultan insuficientes para valorar el eventual exceso de la normal tolerancia. La sentencia tuvo por corroborada la existencia de trabajos de construcción en el lindero y la percepción de ruidos por parte de los testigos, y consideró la coincidencia entre los relatos como un elemento probatorio atendible.

\textbf{Contenido de las declaraciones testimoniales.}
\begin{itemize}
    \item \textbf{Primer testigo} % TODO: contenido incompleto o ambiguo en las notas originales (nombre transcripto con dudas: Pablo Benítez)
    (nombre transcripto con dudas: Pablo Benítez): conocía al actor desde hacía siete años y podía observar la hora al salir y entrar en su domicilio. Declaró que los trabajos en general comenzaban entre las 8 y las 9 de la mañana y se extendían hasta las 10 u 11 de la noche, incluso durante los fines de semana, y que lo sabía porque escuchaba los ruidos de las máquinas y de los obreros. Manifestó que había bastante ruido, constante y muy fuerte, que lo molestaba incluso dentro de su propia vivienda y que en ocasiones vibraban las paredes.
    \item \textbf{Segundo testigo} % TODO: contenido incompleto o ambiguo en las notas originales (nombre no identificado en la grabación)
    (nombre no identificado): durante el año 2024 observó personalmente el inicio de los trabajos y veía diariamente el ingreso de obreros, materiales y maquinaria. Podía escuchar los trabajos desde su domicilio; generalmente comenzaban por la mañana y en varias oportunidades continuaban fuera de los horarios habituales de descanso. Afirmó que los ruidos se escuchaban claramente y eran intensos y constantes.
    \item \textbf{Miriam Lucero}: durante el año 2024 observó los trabajos, el movimiento de obreros, materiales y maquinaria. Explicó que desde su vivienda se podía escuchar la obra, cuyos trabajos comenzaban por la mañana y en varias oportunidades se prolongaban gran parte del día. Sobre los ruidos, afirmó que eran intensos y constantes, que se escuchaban claramente desde su domicilio y que resultaban muy molestos.
\end{itemize}

\textbf{Lectura conjunta.} De estas declaraciones surgen los elementos concretos que debían ponderarse para determinar la normal tolerancia: la intensidad con la que los ruidos eran percibidos desde viviendas vecinas, su persistencia, la extensión temporal de los trabajos y la existencia de jornadas que, según los testimonios, se prolongaban más allá de los horarios habituales de descanso. Es cierto que las declaraciones no constituyen mediciones acústicas ni permiten establecer una cantidad determinada de decibeles, ni cabe atribuirles conclusiones técnicas sobre el nivel de ruido o sobre eventuales daños en la vivienda del actor. Sin embargo, aportan información obtenida mediante la percepción directa de los testigos sobre las características de la actividad y de las molestias en el entorno inmediato. En particular, la declaración del primer testigo aporta una permanencia temporal concreta de los trabajos, desde las 8 hasta las 10 u 11 de la noche, incluso en fines de semana, fuera de los horarios habituales.
% TODO: contenido incompleto o ambiguo en las notas originales (pasaje final de este párrafo ininteligible en la grabación)

Estas circunstancias debían examinarse en conjunto: la actividad desarrollada, las condiciones del lugar y la incidencia de los trabajos sobre los inmuebles vecinos. No se sostiene que cualquier ruido derivado de una construcción constituya una inmisión jurídicamente relevante, sino que se debe determinar si, por su intensidad, duración y horario, las molestias superaron aquello que razonablemente debe tolerarse en las relaciones de vecindad.
\end{example}

\begin{formula}{Art.~1973 del Código Civil y Comercial (según se lo cita en el memorial)}
Las molestias ocasionadas por ruidos, vibraciones y otras emisiones provenientes de inmuebles vecinos no deben exceder la normal tolerancia, atendiendo a las condiciones del lugar y a las circunstancias del caso, aun cuando exista autorización para la actividad.
\end{formula}

\begin{example}{Continuación y cierre del primer agravio}
La norma no establece una tarifa legal de prueba ni exige que el exceso de la normal tolerancia deba acreditarse exclusivamente mediante mediciones acústicas. Por el contrario, la determinación de ese estándar requiere una evaluación circunstanciada de las características de las molestias y de las condiciones concretas en que se producen.

En el escrito se citan además dos artículos del CPCCN: el \textbf{artículo 386}, que faculta al juzgador a apreciar la prueba conforme a las reglas de la sana crítica, y el \textbf{artículo 377}, que establece la carga de acreditar los presupuestos respectivos (síntesis tal como se la formula en la lectura del memorial; el texto de los artículos no fue leído en clase).

En consecuencia, no se pretende desconocer la carga probatoria que recae sobre la parte ni sustituir la prueba objetiva por afirmaciones subjetivas. Lo que se cuestiona es que la sentencia, luego de reconocer la existencia de testimonios coincidentes sobre hechos relevantes, no explica suficientemente por qué las circunstancias descriptas por los declarantes no permiten tener por acreditado el exceso alegado. La ausencia de mediciones objetivas es un elemento que puede ponderarse al resolver, pero no exime de analizar el contenido concreto de las declaraciones ni permite prescindir de las referencias que estas aportan sobre intensidad, horario y duración de los trabajos. La insuficiente consideración de tales circunstancias perjudica a la parte, al conducir a descartar la pretensión fundada en las molestias sin un examen suficiente de los elementos que permitan determinar si estas superaron la normal tolerancia del artículo 1973. Por ello se solicita al tribunal que revise la valoración efectuada en la instancia de grado y determine, a partir del análisis integral de la prueba producida, si las molestias ocasionadas por la obra lindera superaron la tolerancia que razonablemente debe existir.
\end{example}

\section{Responsabilidad civil y daño: no dar nada por supuesto}

\subsection{Comentarios sobre el primer agravio}

El primer agravio contiene una crítica fundamentalmente sobre la \textbf{valoración de la prueba}. Hay prueba: el juzgado reconoce que los testigos declararon y concordaron en que hubo una actividad excesiva. La observación es que el escrito habla poco del daño: se refiere a los hechos y a que la valoración no habría sido adecuada, pero conviene avanzar y sostener que están acreditados los elementos de la responsabilidad civil, porque si se trabaja en torno a la responsabilidad civil hay que demostrar que existe responsabilidad por ese hecho.

La crítica es certera porque va a un punto \textbf{dirimente}, es decir, decisivo, en la médula o el núcleo de la cuestión. Pero hay que llevar al tribunal a la convicción de que esa falencia de la sentencia impidió ver que procede la pretensión y que hay una responsabilidad civil afincada en determinados elementos.

\begin{formula}{Elementos de la responsabilidad civil (según el docente)}
\begin{enumerate}
    \item Daño.
    \item Nexo de causalidad.
    \item Antijuridicidad de la conducta reprochada.
    \item Factor de atribución.
\end{enumerate}
\end{formula}

Es necesario ser categórico en algún punto y volver sobre el daño. De lo contrario, al tribunal le faltaría ese elemento: podría decir que se valoró defectuosamente la prueba, que no se tuvo por acreditado un accionar lesivo del demandado, pero preguntar cuál fue el daño y cómo se mensura. La sentencia también criticaba o negaba la entidad del daño, sobre todo en lo referido a que el actor no habría podido trabajar. Se distinguen el daño material y el extrapatrimonial, sobre los cuales convendría volver en algún punto de la sentencia.

\subsection{Los dos agravios y el daño patrimonial}

El segundo agravio del caso se titula «Examen insuficiente del daño extrapatrimonial». El daño patrimonial no figura en los títulos: el primer agravio, «valoración insuficiente de los ruidos y de la prueba testimonial», lo incluye en cierta forma, pero se concluyó que hay que agregar un apartado específico.

\begin{important}
«No den nada por supuesto»: el tribunal, en principio, no suple omisiones de la parte. Conviene ser incluso redundante, hacer hincapié y dedicar un apartado expreso a cada rubro del daño (patrimonial y extrapatrimonial), conectándolo con los elementos de la responsabilidad civil.
\end{important}

\section{Cuadro Cornell del capítulo}

\begin{longtable}{>{\raggedright\arraybackslash}p{0.27\textwidth}>{\raggedright\arraybackslash}p{0.67\textwidth}}
\toprule
\textbf{Preguntas / palabras clave} & \textbf{Notas / respuestas} \\
\midrule
\endfirsthead
\toprule
\textbf{Preguntas / palabras clave} & \textbf{Notas / respuestas} \\
\midrule
\endhead

\textbf{¿Qué exige la crítica concreta y razonada (art.~265)?} &
La expresión de agravios debe señalar qué parte del fallo es equivocada y por qué. No alcanza con decir «se violó la defensa en juicio»: hay que mostrar concretamente cómo y por qué una decisión lesiona un precepto determinado. \\[0.6em]

\textbf{¿Qué debe expresar el objeto del escrito?} &
El propósito del escrito, que se desarrollará en los apartados siguientes. Adelantar peticiones propias del petitorio no está mal, pero es una cuestión de forma. \\[0.6em]

\textbf{¿Cómo se estructuran los agravios?} &
Con encabezados como «primer agravio» y «segundo agravio», o directamente con el nombre de cada agravio. En el caso: valoración insuficiente de los ruidos y de la prueba testimonial, y examen insuficiente del daño extrapatrimonial; luego costas y honorarios. \\[0.6em]

\textbf{¿Qué significa que el recurso sea autosuficiente y por qué se parafrasea la sentencia?} &
Su lectura debe permitir entender toda la causa. Se reconstruye primero qué decidió el juez y cómo (paráfrasis lo más textual posible, identificando el considerando) y recién después se critica. \\[0.6em]

\textbf{¿Cómo se critica la valoración de la prueba?} &
Reconociendo lo que la sentencia admitió (testimonios coincidentes), mostrando que no explicó por qué no alcanzaban y apoyándose en las normas aplicables: art.~1973 CCyC (normal tolerancia), sana crítica (art.~386) y carga de la prueba (art.~377). \\[0.6em]

\textbf{¿Qué hace que un agravio sea «dirimente»?} &
Que toque el núcleo de la cuestión y lleve al tribunal a la convicción de que están acreditados los elementos de la responsabilidad civil: daño, nexo de causalidad, antijuridicidad y factor de atribución. \\[0.6em]

\textbf{¿Qué significa «no den nada por supuesto»?} &
El tribunal no suple omisiones. Hay que ser incluso redundante y dedicar un apartado expreso al daño patrimonial y al extrapatrimonial. \\

\midrule
\multicolumn{2}{>{\raggedright\arraybackslash}p{0.94\textwidth}}{%
\textbf{Resumen:} recurrir no es repetir el reclamo. La expresión de agravios debe tener un objeto claro, agravios ordenados y titulados, ser autosuficiente y formular una crítica concreta y razonada que reconstruya lo decidido y muestre por qué es equivocado y con qué norma. Para que la crítica sea dirimente debe conectarse con los elementos de la responsabilidad civil, sin dar nada por supuesto.
} \\
\bottomrule
\end{longtable}

% ----------------------------------------------------------
\chapter[Prueba en segunda instancia y trámite]{Prueba en segunda instancia y trámite del recurso}
% ----------------------------------------------------------

\section{Apertura a prueba en segunda instancia (artículo 260)}

\subsection{Planteo}

Se reflexionó sobre la posibilidad de abrir a prueba en la segunda instancia. Aunque es muy restrictiva, no debe desmerecerse: si hubo una hipótesis de \textbf{negligencia}, puede haber prueba para replantear. En el caso práctico no correspondía, porque no se lo planteó así en el caso ficcionado.

\subsection{Caso real: el apelante que acompaña documental con el memorial}

\begin{example}{Apelante que acompaña documental con el memorial}
En un recurso que debió contestarse, la parte apelante acompañó una documental con el memorial de apelación. Esos documentos habían estado sometidos al conocimiento del juez de primera instancia, pero se acompañaron nuevamente. La parte apelada desconoció los documentos y sostuvo que no era la oportunidad, que no correspondía y que la vía no es para acompañar nueva prueba, tanto más tratándose de un recurso de apelación \emph{en relación}. Si se tratara de un recurso libre, podría llegarse por el procedimiento del art.~260. Se buscó además reivindicar que hay una cosa ya resuelta, una cosa juzgada, respecto de la parte que se intenta apelar.
\end{example}

Los apelantes intentan muchas veces por todos los medios atacar la resolución e incorporan elementos que no estaban sometidos al conocimiento del juez de primera instancia. El recurso de apelación \textbf{no es un nuevo juicio}: es un reexamen de lo que ya se decidió.

\begin{important}
El recurso de apelación no es un nuevo juicio, sino un reexamen de lo ya decidido. Acompañar prueba nueva con el memorial, sobre todo en una apelación en relación, no corresponde: ese material debió articularse en primera instancia.
\end{important}

\subsection{Negligencia, denegación de prueba y artículo 260}

Frente a la consulta sobre si puede replantearse en segunda instancia una prueba pericial que no se produjo, la respuesta fue que habría que ver las circunstancias del caso. El artículo 260 (siempre en el recurso de apelación libre, no en el recurso en relación) contempla la posibilidad de replantear pruebas en los supuestos de negligencia o de denegación.

Cuando el juez de primera instancia deniega una medida de prueba, no puede apelarse: hay restricción del recurso de apelación. Entonces la cámara puede examinar si el juez denegó mal la medida. Esto parece un dispendio de actividad jurisdiccional, porque hay que transitar todo el juicio con esa denegatoria y recién después se tiene la oportunidad de plantearlo ante la cámara, que puede decir que corresponde hacer la pericia. En tal caso habrá una etapa probatoria bastante circunscripta en la cámara.

Como ejemplo de denegación, a veces se pide pericia médica y psicológica y el juzgado dispone que la psicológica la haga el mismo médico legista, que luego no puede profundizarla como correspondería; la parte se considera vulnerable y puede replantearlo en la cámara. Es una opción teórica.

En el caso práctico, como se hicieron cuatro o cinco pasos procesales, no se profundizó en estos temas, y la sentencia se apoyó sobre todo en que no había pericia: el expediente llegó sin ella.

Habría que leer con precisión el artículo 260, en toda su dimensión (incluso «en cada inciso»), para ver la posibilidad de replantear prueba. Muy pocos abogados se animan a hacerlo. Generalmente, el recurso de apelación libre versa sobre lo ya actuado y no se agrega nada nuevo.

% TODO: contenido incompleto o ambiguo en las notas originales (el texto del art. 260 no fue leído en clase; no se precisa si los supuestos son solo negligencia o denegación)
\begin{note}
El texto del artículo 260 del CPCCN no fue leído en clase; el docente lo cita como la vía para el replanteo de prueba en segunda instancia en la apelación libre y recomienda leerlo en toda su dimensión.
\end{note}

\section{Recursos con efecto diferido y restricción de apelar en la etapa de prueba}

Los recursos concedidos con \textbf{efecto diferido} deben hacerse valer nuevamente en la oportunidad correspondiente: si no se vuelve sobre ellos en esa oportunidad, se consideran \textbf{desiertos}. Hay casos en que no hubo posibilidad de apelar porque la normativa no lo permite: la restricción en materia de apelación en la etapa de prueba es intensa. Se menciona el artículo 379, que niega la posibilidad de apelar.

% TODO: contenido incompleto o ambiguo en las notas originales (el docente cita el art. 379 con dudas: «creo, si no me equivoco»; su texto no fue leído)
\begin{note}
Según el docente, el artículo 379 del CPCCN niega la posibilidad de apelar en la etapa de prueba (lo cita con reservas: «creo, si no me equivoco»). Su texto no fue leído en clase.
\end{note}

\section{Presentación del recurso: firma y plataformas}

\subsection{La firma del escrito}

Un escrito sin firma es un \textbf{escrito inexistente}. Cuando se actúa por derecho propio hay que acompañar el escrito firmado, y es un olvido frecuente. Debe conocerse además el procedimiento de cada jurisdicción: en la jurisdicción nacional se sube directamente el PDF firmado, mientras que en la provincia de Buenos Aires hay que pegar el texto en una plantilla y subir el escrito como adjunto firmado, lo que implica un doble paso. Si el escrito fue elaborado por varias personas, basta con que tenga una firma.

\begin{important}
Un escrito sin firma es un escrito inexistente. Hay que verificar siempre la firma antes de presentarlo y conocer el procedimiento de cada jurisdicción: en la jurisdicción nacional se sube el PDF firmado; en la provincia de Buenos Aires se pega el texto en una plantilla y se adjunta el escrito firmado (doble paso).
\end{important}

\section{Cuadro Cornell del capítulo}

\begin{longtable}{>{\raggedright\arraybackslash}p{0.27\textwidth}>{\raggedright\arraybackslash}p{0.67\textwidth}}
\toprule
\textbf{Preguntas / palabras clave} & \textbf{Notas / respuestas} \\
\midrule
\endfirsthead
\toprule
\textbf{Preguntas / palabras clave} & \textbf{Notas / respuestas} \\
\midrule
\endhead

\textbf{¿Puede producirse prueba en segunda instancia?} &
Solo excepcionalmente, en la apelación libre y en los supuestos del art.~260 (por ejemplo, replanteo por negligencia o denegación de la medida). Muy pocos abogados lo hacen; generalmente la apelación libre versa sobre lo ya actuado. \\[0.6em]

\textbf{¿Por qué la apelación no es un nuevo juicio?} &
Porque es un reexamen de lo ya decidido. Acompañar documental nueva con el memorial, sobre todo en una apelación en relación, no corresponde: debió articularse en primera instancia. \\[0.6em]

\textbf{¿Qué pasa cuando el juez de primera instancia deniega una prueba?} &
La denegatoria no es apelable (restricción del recurso). La cámara puede examinar si el juez denegó mal la medida, lo que implicaría una etapa probatoria circunscripta en la cámara. \\[0.6em]

\textbf{¿Qué ocurre con un recurso concedido con efecto diferido?} &
Debe hacerse valer nuevamente en la oportunidad correspondiente; de lo contrario se declara desierto. En la etapa de prueba la apelación está muy restringida (art.~379, según el docente). \\[0.6em]

\textbf{¿Qué ocurre si se presenta un escrito sin firma?} &
Es un escrito inexistente. Hay que verificar la firma y el procedimiento de cada jurisdicción (PDF firmado en la jurisdicción nacional; plantilla más adjunto firmado en provincia de Buenos Aires). \\

\midrule
\multicolumn{2}{>{\raggedright\arraybackslash}p{0.94\textwidth}}{%
\textbf{Resumen:} la apelación no es un nuevo juicio y la prueba en segunda instancia es excepcional (art.~260, apelación libre). Los recursos concedidos con efecto diferido deben retomarse en la oportunidad correspondiente o se declaran desiertos, y el escrito debe estar firmado para existir.
} \\
\bottomrule
\end{longtable}

% ----------------------------------------------------------
\chapter[Apelación de honorarios]{La apelación de honorarios}
% ----------------------------------------------------------

\section{Honorarios, costas y autonomía}

\subsection{Presentación del tema}

Cómo se apelan los honorarios es un tema que forma parte de la cotidianeidad del ejercicio profesional. En el caso práctico no se tuvo esa posibilidad, porque la sentencia difirió la regulación de honorarios; pero hay sentencias que tienen la regulación de honorarios contenida.

\subsection{Costas, honorarios y su relación con el éxito de la demanda}

La sentencia decide sobre un tema principal y sobre accesorios. Los accesorios vienen de la mano de las costas: la sentencia «impone las costas y regula los honorarios». Un recurso libre contra la sentencia ataca la parte principal, la imposición de costas y también la regulación de honorarios.

Si se logra que la sentencia se revoque y se haga lugar a la demanda, cambia la imposición de costas: las costas siguen la suerte de lo principal, es decir, se imponen al vencido, que en tal caso sería el demandado (en primera instancia el actor fue el vencido). Los honorarios tienen cierta \textbf{autonomía}: fueron regulados más allá de la imposición de costas. Su contenido se relaciona con el progreso o el rechazo de la demanda, porque si se gana el juicio el honorario para la parte vencedora es mayor. Siempre habrá un honorario, salvo casos excepcionales en que se considera que el trabajo es inoficioso.

\subsection{El trabajo inoficioso y el pacto de cuota litis}

A veces los jueces consideran que el trabajo profesional es \emph{inoficioso}, como si no tuviera valor. Los profesionales suelen estar ligados al éxito, es decir, a que prospere la demanda. Se relató un caso en que no prosperó la demanda y existía un \textbf{pacto de cuota litis} ligado al éxito (un 20\,\% si prosperaba la demanda): solo se cobró un gasto, que luego se rindió.

\subsection{Regulación autónoma o contenida en la sentencia}

Puede haber una regulación de honorarios en una resolución autónoma y otra contenida en otro decisorio. Si el juzgado dicta un auto regulatorio autónomo, o si la sentencia difiere la regulación de honorarios y luego los regula, la regulación tiene un tratamiento autónomo.

\begin{table}[H]
\centering
\begin{tabularx}{\textwidth}{>{\raggedright\arraybackslash}p{0.2\textwidth}>{\raggedright\arraybackslash}X>{\raggedright\arraybackslash}X}
\toprule
\textbf{Aspecto} & \textbf{Regulación autónoma (auto regulatorio)} & \textbf{Regulación contenida en la sentencia} \\
\midrule
\textbf{Dónde está} & En una resolución separada, dictada luego de la sentencia (por ejemplo, cuando la sentencia difiere la regulación). & Dentro de la propia sentencia, junto con la imposición de costas. \\
\textbf{Vía de ataque} & Apelación específica de la regulación, con régimen propio del artículo 244. & Se discute si se apela junto con la sentencia o si exige apelación específica (ver sección sobre honorarios contenidos en la sentencia). \\
\textbf{Relación con las costas} & Autonomía: la regulación existe más allá de la imposición de costas. & Depende en parte del resultado del recurso: si se revoca la imposición de costas, se modifica la regulación. \\
\bottomrule
\end{tabularx}
\end{table}

\section{Artículo 244: «interponerse» y «fundarse»}

Para la apelación de una regulación autónoma el artículo aplicable es el 244 del Código Procesal (el 242 inicia el recurso de apelación).

\begin{formula}{Art.~244 CPCCN (según lectura en clase)}
Habiendo disposición en contrario, el plazo para apelar será de cinco días. \textbf{Toda regulación de honorarios será apelable.} El recurso de apelación deberá interponerse y podrá fundarse dentro de los cinco días de la notificación.
\end{formula}

% TODO: contenido incompleto o ambiguo en las notas originales (la fórmula inicial se transcribe tal como fue leída por un alumno)

La expresión «toda regulación de honorarios será apelable» es un precepto muy importante y garantista: el artículo 242 restringe el recurso de apelación de acuerdo con un monto mínimo, y si no se trasciende ese mínimo no se puede apelar. Con el artículo 244, en la jurisdicción nacional, la regulación puede apelarse aunque el monto sea bajo, dentro de los cinco días.

\subsection{Dos verbos modales: «deberá» y «podrá»}

El artículo emplea dos verbos modales: \textbf{«deberá»}, que indica obligatoriedad, y \textbf{«podrá»}, que indica facultad.

\begin{table}[H]
\centering
\begin{tabularx}{\textwidth}{>{\raggedright\arraybackslash}p{0.25\textwidth}>{\raggedright\arraybackslash}X>{\raggedright\arraybackslash}p{0.2\textwidth}}
\toprule
\textbf{Verbo del artículo 244} & \textbf{Acto procesal} & \textbf{Carácter} \\
\midrule
\textbf{«Deberá» interponerse} & Interponer el recurso de apelación dentro de los 5 días de la notificación. & Obligatorio \\
\textbf{«Podrá» fundarse} & Fundar el recurso (expresar los motivos) dentro de los mismos 5 días. & Facultativo \\
\bottomrule
\end{tabularx}
\end{table}

\subsection{Apelar «por altos» o «por bajos» sin fundar}

Con frecuencia se dice «apelo» y «por bajos» o «por altos», sin ningún fundamento. Por ejemplo, ante una regulación de diez UMA, quien debe pagar los honorarios por la imposición de costas los apela por altos, y con ello pide que la cámara revise si son adecuados. No es imperativo fundar: basta decir «apelo», aunque se debe indicar mínimamente el sentido («por altos» o «por bajos»). Aunque no se expresen los motivos, ello amerita el examen de la cámara, que revisará si la regulación ha sido adecuada y puede dar la razón o confirmar. Quien es beneficiario de los honorarios apela por bajos si considera que le corresponde más.

\section{Cómo fundar la apelación de honorarios y la Ley de Aranceles}

\subsection{Modelos de agravios}

Al apelar honorarios, los agravios pueden expresarse, por ejemplo, de las siguientes maneras:

\begin{itemize}
    \item Insuficiente ponderación de la labor profesional desarrollada.
    \item Inadecuada valoración de la calidad jurídica de la labor.
    \item Falta de consideración de la complejidad del asunto.
    \item Omisión de ponderar la extensión temporal de la intervención profesional.
    \item Inadecuada determinación de la base regulatoria.
    \item Apartamiento de las pautas legales de regulación.
    \item Insuficiente consideración del resultado obtenido.
    \item Falta de reconocimiento de las etapas procesales cumplidas.
    \item Insuficiente ponderación de la responsabilidad profesional asumida.
    \item Falta de proporcionalidad entre la retribución y el trabajo realizado.
    \item Insuficiente fundamentación de la resolución regulatoria.
\end{itemize}

Todas esas cuestiones están en la \textbf{Ley de Aranceles}, a la que hay que acudir para dar contenido a la apelación si se desea fundarla; se mencionó fundamentalmente el artículo 16 de dicha ley. La base está en el artículo 244 del Código Procesal.

Esto es desafiante porque, a diferencia de las resoluciones intermedias del proceso que se recurren en relación (apelo; el juzgado concede o deniega; se notifica la concesión por ministerio de la ley; a partir de allí se funda, y si no se funda hay deserción), aquí quien quiera fundar debe hacerlo \emph{en el mismo lapso} en que apela: dentro de los mismos 5 días.

% TODO: contenido incompleto o ambiguo en las notas originales (el artículo de la Ley de Aranceles se menciona «creo que fundamentalmente aludió al artículo 16»)

\subsection{Ventajas y riesgos de fundar}

Es de buena práctica fundar. Cuando se dice «apelo por bajos» o «apelo por altos» sin más, la otra parte no puede decir nada respecto de esa apelación; solo podría apelar a su vez, por altos o por bajos. En cambio, si se argumenta, aunque sea con una palabra más, el juzgado dispone «de los fundamentos, traslado», y la otra parte puede contestar que el honorario está perfecto, que hay que reducirlo o que hay que aumentarlo. Se le abre la puerta a la otra parte. Esto no debe intimidar, pero deben pensarse las consecuencias: quien fundó ejerció la facultad de fundamentar y se genera ese circuito.

\begin{table}[H]
\centering
\begin{tabularx}{\textwidth}{>{\raggedright\arraybackslash}p{0.2\textwidth}>{\raggedright\arraybackslash}X>{\raggedright\arraybackslash}X}
\toprule
\textbf{Opción} & \textbf{Apelar sin fundar («apelo por altos / por bajos»)} & \textbf{Apelar y fundar} \\
\midrule
\textbf{Exigencia legal} & Basta con la palabra «apelo»; fundar es facultativo (art.~244). & Se ejerce la facultad de fundar dentro de los mismos 5 días. \\
\textbf{Efecto en el trámite} & La cámara igualmente revisa si la regulación fue adecuada; la contraparte solo puede apelar a su vez. & El juzgado corre traslado de los fundamentos; la contraparte puede contestar y abrir el debate. \\
\textbf{Riesgo / ventaja} & Menor exposición; la revisión se produce igual. & Buena práctica y mayor contenido, pero se «abre la puerta» a la otra parte. \\
\bottomrule
\end{tabularx}
\end{table}

Todo lo anterior se refiere a un auto regulatorio autónomo.

\section{Honorarios contenidos en la sentencia: ¿cuándo fundar?}

\subsection{El planteo}

Cuando se apela la sentencia en un proceso ordinario, el juzgado de primera instancia concede el recurso (salvo que no corresponda, por ejemplo por no superar el umbral de apelabilidad del art.~242) y el expediente va a la cámara, donde se abre el espectro de los arts.~259 y 260. En el proceso sumarísimo la apelación es en relación.

La cuestión es qué ocurre con los honorarios contenidos en la sentencia: si van con esa apelación o si hay que hacer una apelación específica. Se señaló que en el fuero laboral se apela la sentencia y también los honorarios, porque se apela y se funda; en civil está desdoblado.

Quien apeló la sentencia puede decir «apelo la sentencia y apelo los honorarios por altos» (o «por bajos»), sin decir más. Pero si es beneficiario de los honorarios (por ejemplo, la parte demandada que logró el rechazo de la demanda) y considera insuficientes los honorarios fijados, debería sumar fundamentos por mandato del artículo 244, que se refiere específicamente a los honorarios.

\subsection{El riesgo de la extemporaneidad}

Si no se fundan los honorarios al apelar, la cámara tratará el fondo del asunto. Si confirma la sentencia, y la parte había apelado los honorarios «por bajos» y acompaña en la expresión de agravios un argumento sobre el \emph{quantum} de los honorarios, se le dirá que es \textbf{extemporáneo}: debió hacerlo en la oportunidad del artículo 244. La normativa no contempla fundar los honorarios junto con los agravios de la sentencia: da una directiva específica de apelación de honorarios (dentro de los 5 días se pueden apelar y fundar). Los 5 días coinciden con el plazo para apelar la sentencia, pero se desdobla la fundamentación respecto de los honorarios.

La práctica sugerida es: si hay argumentos sobre los honorarios, expresarlos en el momento de apelar, dentro de los 5 días, indicando por qué son bajos o altos. Si no se tienen motivos, se pone toda la energía en que la sentencia se revoque y se cambie la imposición de costas, lo cual puede modificar por añadidura la regulación de honorarios de primera instancia.

\begin{example}{Costas y regulación de honorarios (hipótesis)}
Se rechazó la demanda, se impusieron las costas al actor y se regularon honorarios: 10 UMA para el abogado del actor (cuya demanda fue rechazada) y 20 UMA para el abogado del demandado. Se puede apelar para que se revoque la sentencia y la imposición de costas, lo cual traerá por añadidura la regulación de honorarios. Pero si se considera que los honorarios fijados a favor del abogado de la demandada son altos y hay razones para decirlo, debe decirse al apelar, por la hipótesis de que se confirme la sentencia: si se confirma la sentencia, se confirman los honorarios.
\end{example}

Apelar los honorarios por altos o por bajos sirve también para que no queden firmes y vayan a la cámara, lo que da al cliente un poco más de tiempo para pagar, porque los honorarios no están firmes.

\begin{important}
Cuando los honorarios están contenidos en la sentencia, conviene fundar la apelación de honorarios en el mismo acto de apelar (dentro de los 5 días), si hay motivos. Si se guardan los argumentos para la expresión de agravios, la cámara puede considerarlos extemporáneos (art.~244). Apelar la imposición de costas puede, por añadidura, modificar la regulación, pero no la corrige si la sentencia se confirma.
\end{important}

\section{Apelación libre y apelación en relación: ¿cuándo se funda?}

Cuando el recurso se concede \textbf{libremente}, se funda en la cámara; en el proceso ordinario, el plazo para fundarlo es de 10 días. Cuando es \textbf{en relación} (por ejemplo, la sentencia de juicio sumarísimo), se apela en 5 días; cuando se concede el recurso, la concesión se notifica por ministerio de la ley y a partir de allí corre el plazo para presentar el memorial, que se presenta en primera instancia. Hoy en día, la apelación libre corresponde solamente a la sentencia de juicio ordinario.

\begin{example}{Cómputo del memorial en la apelación en relación}
Si la concesión del recurso se dictó hoy, se notifica por ministerio de la ley el viernes y se tiene toda la semana siguiente para presentar el memorial. Como el lunes 23 es feriado, se pasa a las dos primeras horas del martes.
\end{example}

\begin{table}[H]
\centering
\begin{tabularx}{\textwidth}{>{\raggedright\arraybackslash}p{0.2\textwidth}>{\raggedright\arraybackslash}X>{\raggedright\arraybackslash}X}
\toprule
\textbf{Aspecto} & \textbf{Apelación libre} & \textbf{Apelación en relación} \\
\midrule
\textbf{Ejemplo} & Sentencia definitiva de proceso ordinario. & Sentencia de juicio sumarísimo; resoluciones intermedias del proceso. \\
\textbf{Dónde se funda} & En la cámara (segunda instancia). & En primera instancia, mediante memorial. \\
\textbf{Plazo para fundar} & 10 días (proceso ordinario). & Desde la notificación de la concesión por ministerio de la ley; en el ejemplo, la semana siguiente. \\
\textbf{Prueba en segunda instancia} & Posible en los supuestos del artículo 260. & No corresponde acompañar prueba nueva. \\
\bottomrule
\end{tabularx}
\end{table}

\section{Ley de Aranceles 27.423: notificación, pago, mora e intereses}

La \textbf{Ley 27.423} es una herramienta que debe tenerse siempre presente para defender los honorarios. La ley anterior, la 21.839, fijaba montos (por ejemplo, de cinco mil pesos) que quedaban sin actualizarse, y a lo sumo sumaban tasa de interés.

\subsection{Artículo 57: cobro al cliente}

\begin{formula}{Art.~57 Ley 27.423 (según lectura en clase)}
Si el condenado en costas no abonara los honorarios profesionales en tiempo y forma, el profesional podrá requerir su pago a su cliente luego de 30 días corridos de intimado, y siempre que esté debidamente notificado.
\end{formula}

Según el comentario del docente, el cliente queda obligado: en principio se reclama al condenado en costas, pero el profesional puede reclamar a su propio cliente si este no es el condenado en costas.

\subsection{Artículo 56: notificación y apelación}

\begin{formula}{Art.~56 Ley 27.423 (según lectura en clase)}
Las resoluciones que regulan honorarios deberán ser notificadas a sus beneficiarios y a los obligados al pago personalmente o por cédula, o por cualquier medio previsto en el Código Procesal Civil y Comercial de la Nación. [Podrán apelarse] en el término de 5 días, pudiendo fundarse la apelación en el acto de interponerse el recurso. En caso de apelación de honorarios, será de aplicación lo dispuesto en el artículo 244 del Código Procesal Civil y Comercial de la Nación. Las cámaras de apelaciones deberán resolver el recurso dentro de los 30 días de recibido el expediente. Los honorarios serán apelables con prescindencia de los montos de los mismos. En el caso de que la regulación forme parte de la sentencia definitiva, en la notificación se deberá acompañar copia íntegra de la misma, bajo pena de nulidad de la notificación.
\end{formula}

Esta norma se conecta con el artículo 244, que es la norma fundante. La exigencia de acompañar la copia íntegra de la sentencia \textbf{bajo pena de nulidad} busca garantizar al interesado el conocimiento íntegro. Hoy, en general, la notificación electrónica va con el adjunto completo.

% TODO: contenido incompleto o ambiguo en las notas originales (el docente indica que falta mencionar otro artículo de la Ley de Aranceles sobre notificación, cuyo número no recuerda, y propone buscarlo para la próxima clase)

\subsection{Honorarios de segunda instancia}

Un alumno leyó una disposición sobre los honorarios de segunda instancia: «Si la sentencia recurrida fuera revocada en todas sus partes en favor del apelante, los honorarios profesionales por los trabajos en esa instancia de apelación se fijarán entre el 30 y el 40\,\% de los correspondientes a la primera instancia». El docente aclaró que esa disposición trata de los honorarios de apelación, mientras que el tema en discusión era qué ocurre si se confirma la sentencia, en particular con la mora.

% TODO: contenido incompleto o ambiguo en las notas originales (no se indica el número de artículo de la disposición leída por el alumno sobre honorarios de segunda instancia)

\subsection{Artículo 54: plazo de pago, ejecución e intereses}

Si se apeló por bajos o por altos y la cámara confirma, el plazo para pagar los honorarios surge del artículo 54: \textbf{10 días} (antes eran 30; la nueva ley abrevió ese plazo).

\begin{formula}{Art.~54 Ley 27.423 (según lectura en clase)}
Los honorarios regulados judicialmente deberán abonarse dentro de los diez días de quedar firme la resolución regulatoria. Los honorarios extrajudiciales, dentro de los diez días de intimado su pago, cuando sean exigibles. Los honorarios a cargo del mandante o patrocinado deberán serle notificados al domicilio real o al constituido al efecto. La acción por cobro de honorarios regulados judicialmente tramitará por la vía de ejecución de sentencia. En ningún caso abonará tasa de justicia ni estará sujeta a ningún tipo de contribución. Los honorarios deberán pagarse siempre en moneda de curso legal. Las deudas de honorarios pactados o por regulación judicial firme, cuando hubiere mora del deudor, devengarán intereses desde la fecha de la regulación de primera instancia y hasta el momento de su efectivo pago, los que serán fijados por el juez de la causa, siguiendo el mismo criterio que el utilizado para establecer la actualización de los valores económicos de la causa.
\end{formula}

\subsection{¿Cuándo queda firme la regulación?}

La regulación de primera instancia queda firme a los cinco días. Si es de cámara, queda firme a los 10 días, porque puede ser susceptible de recursos extraordinarios; por lo tanto hay prácticamente 20 días.

\subsection{Apelar honorarios y el riesgo de incurrir en mora}

\begin{example}{Mora e intereses retroactivos}
Se apelan honorarios por altos o por bajos en noviembre de 2026 y la cámara se expide en marzo de 2027 (posibilidad planteada como hipótesis). La cámara los confirma, sin bajarlos ni subirlos. El obligado debe estar atento para pagarlos dentro de los 10 días, porque de lo contrario incurre en mora y adeuda intereses retroactivos a la regulación de primera instancia. Por eso, apelar por bajos o por altos debe considerar, para el obligado al pago, este aspecto: no incurrir en mora.
\end{example}

Respecto de la tasa de interés, el docente ejemplificó con una hipótesis: si la inflación fuera del 2\,\% mensual y la tasa de interés del 1\,\% mensual (y el plazo fijo diera como máximo un 20\,\% anual frente a una inflación del 24\,\%), el diferimiento no afectaría tanto. Aun así, si se paga dentro de esos 10 días se evita pagar intereses.

Esto aplica \textbf{solo si la cámara confirma}: si los aumenta o los baja, sería una nueva regulación, dictada directamente por la cámara.

\subsection{Cuadro de plazos y reglas}

\begin{table}[H]
\centering
\begin{tabularx}{\textwidth}{>{\raggedright\arraybackslash}p{0.27\textwidth}>{\raggedright\arraybackslash}p{0.22\textwidth}>{\raggedright\arraybackslash}X}
\toprule
\textbf{Plazo / regla} & \textbf{Norma} & \textbf{Contenido} \\
\midrule
5 días para apelar (y fundar, si se quiere) & CPCCN art.~244; Ley 27.423 art.~56 & Toda regulación es apelable; «deberá» interponerse y «podrá» fundarse en el mismo plazo. \\
Apelación sin importar el monto & Ley 27.423 art.~56 & Los honorarios son apelables con prescindencia de su monto. \\
30 días para que la cámara resuelva & Ley 27.423 art.~56 & Plazo de las cámaras desde que reciben el expediente. \\
Copia íntegra de la sentencia en la notificación & Ley 27.423 art.~56 & Si la regulación forma parte de la sentencia definitiva: bajo pena de nulidad de la notificación. \\
10 días para pagar desde que queda firme & Ley 27.423 art.~54 & Antes eran 30 días. En primera instancia, firmeza a los 5 días; en cámara, a los 10 por posibles recursos extraordinarios ($\approx$ 20 días). \\
Intereses por mora & Ley 27.423 art.~54 & Desde la regulación de primera instancia hasta el pago efectivo, si hubo mora. \\
Cobro al cliente tras 30 días corridos de intimación & Ley 27.423 art.~57 & Si el condenado en costas no paga, el profesional puede reclamar a su cliente. \\
Honorarios de segunda instancia: 30\,\% a 40\,\% de los de primera & Ley 27.423 (artículo leído por un alumno) & Cuando la sentencia es revocada en todas sus partes a favor del apelante. \\
\bottomrule
\end{tabularx}
\end{table}

\section{Cuadro Cornell del capítulo}

\begin{longtable}{>{\raggedright\arraybackslash}p{0.27\textwidth}>{\raggedright\arraybackslash}p{0.67\textwidth}}
\toprule
\textbf{Preguntas / palabras clave} & \textbf{Notas / respuestas} \\
\midrule
\endfirsthead
\toprule
\textbf{Preguntas / palabras clave} & \textbf{Notas / respuestas} \\
\midrule
\endhead

\textbf{¿Cómo se relacionan costas y honorarios?} &
Las costas siguen la suerte de lo principal; los honorarios tienen cierta autonomía, pero su monto se vincula con el éxito o rechazo de la demanda. La regulación puede ser autónoma o estar contenida en la sentencia. \\[0.6em]

\textbf{Art.~244: ¿qué es obligatorio y qué es facultativo?} &
Toda regulación de honorarios es apelable. El recurso «deberá» interponerse (obligatorio) y «podrá» fundarse (facultativo) dentro de los 5 días de la notificación. \\[0.6em]

\textbf{¿Qué significa apelar «por altos» o «por bajos»?} &
Basta decir «apelo» e indicar el sentido: la cámara igualmente revisa si la regulación fue adecuada, aunque no se expresen los motivos. \\[0.6em]

\textbf{¿Conviene fundar la apelación de honorarios?} &
Es de buena práctica, pero abre traslado a la otra parte, que puede contestar. Se pueden usar los criterios de la Ley de Aranceles (labor, complejidad, extensión, base regulatoria, resultado, proporcionalidad, etc.). \\[0.6em]

\textbf{Honorarios dentro de la sentencia: ¿cuándo fundar?} &
En el mismo acto de apelar, dentro de los 5 días, si hay motivos; de lo contrario, los argumentos presentados luego en la expresión de agravios pueden rechazarse por extemporáneos. \\[0.6em]

\textbf{Apelación libre vs. en relación: ¿dónde y cuándo se funda?} &
Libre (sentencia de proceso ordinario): se funda en cámara, 10 días. En relación (por ejemplo, sumarísimo): memorial en primera instancia tras la notificación por ministerio de la ley. \\[0.6em]

\textbf{Plazos de la Ley 27.423 (arts.~54, 56 y 57)} &
Apelación en 5 días; la cámara resuelve en 30; pago en 10 días desde la firmeza ($\approx$ 20 días en cámara); notificación con copia íntegra bajo pena de nulidad; reclamo al cliente tras 30 días corridos de intimación. \\[0.6em]

\textbf{¿Qué riesgo hay al apelar honorarios?} &
Si se confirman y no se pagan dentro de los 10 días, hay mora: intereses desde la regulación de primera instancia. Si la cámara los modifica, es una nueva regulación. \\

\midrule
\multicolumn{2}{>{\raggedright\arraybackslash}p{0.94\textwidth}}{%
\textbf{Resumen:} el cobro de los honorarios es parte esencial del ejercicio profesional. Toda regulación es apelable en 5 días; fundar es una facultad con ventajas y riesgos (traslado a la otra parte, posible extemporaneidad); la cámara tiene 30 días para resolver y el pago debe hacerse en 10 días desde la firmeza para no incurrir en mora e intereses.
} \\
\bottomrule
\end{longtable}

% ----------------------------------------------------------
\chapter[Cuadro Cornell general]{Cuadro Cornell general de repaso}
% ----------------------------------------------------------

\begin{longtable}{>{\raggedright\arraybackslash}p{0.27\textwidth}>{\raggedright\arraybackslash}p{0.67\textwidth}}
\toprule
\textbf{Preguntas / palabras clave} & \textbf{Notas / respuestas} \\
\midrule
\endfirsthead
\toprule
\textbf{Preguntas / palabras clave} & \textbf{Notas / respuestas} \\
\midrule
\endhead

\textbf{¿Cuándo se abre el recurso extraordinario en un caso de derecho común?} &
Por la doctrina de la arbitrariedad o cuando se declaró la inconstitucionalidad de una norma, que abre la cuestión federal (ley 48, art.~14). \\[0.6em]

\textbf{¿Qué es la crítica concreta y razonada?} &
La exigencia del art.~265: señalar qué parte del fallo es equivocada, por qué y con qué norma, sin quedarse en afirmaciones generales como «se violó la defensa en juicio». \\[0.6em]

\textbf{¿Cómo se construye un buen escrito de agravios?} &
Objeto claro (propósito, sin adelantar todo el petitorio); agravios ordenados y titulados; paráfrasis lo más textual posible de la sentencia, identificando el considerando, y recién después la crítica. El recurso debe ser autosuficiente. \\[0.6em]

\textbf{¿Qué hace dirimente a un agravio y qué significa «no den nada por supuesto»?} &
Tocar el núcleo de la cuestión y conectarlo con los elementos de la responsabilidad civil (daño, nexo de causalidad, antijuridicidad, factor de atribución). El tribunal no suple omisiones: hay que dedicar apartados expresos a cada rubro. \\[0.6em]

\textbf{¿Puede producirse prueba en segunda instancia?} &
Solo excepcionalmente, en la apelación libre y en los supuestos del art.~260. La apelación no es un nuevo juicio. \\[0.6em]

\textbf{¿Qué advertencias prácticas surgen sobre el trámite del recurso?} &
Los recursos concedidos con efecto diferido deben retomarse o se declaran desiertos; un escrito sin firma es inexistente. \\[0.6em]

\textbf{¿Cómo se apelan los honorarios?} &
Toda regulación es apelable en 5 días; el recurso «deberá» interponerse y «podrá» fundarse en ese plazo (art.~244). Si están en la sentencia, conviene fundar al apelar. \\[0.6em]

\textbf{¿Qué plazos y riesgos prevé la Ley 27.423?} &
Pago en 10 días desde la firmeza, cámara resuelve en 30 días, notificación con copia íntegra bajo pena de nulidad, reclamo al cliente tras 30 días corridos de intimación y mora con intereses desde la regulación de primera instancia. \\[0.6em]

\textbf{¿Qué rol tiene la inteligencia artificial?} &
Es una herramienta útil, pero el criterio, la estrategia y la responsabilidad profesional siguen siendo del abogado. \\

\midrule
\multicolumn{2}{>{\raggedright\arraybackslash}p{0.94\textwidth}}{%
\textbf{Síntesis final:} recurrir no es repetir el reclamo, sino \textbf{reconstruir lo que decidió el juez y criticarlo de manera concreta y razonada}, mostrando qué parte es equivocada, por qué y con qué norma, sin dar nada por supuesto. Esa crítica se apoya en una Constitución que se vive y se aplica (arbitrariedad, cuestión federal, derechos humanos) y se mueve dentro de límites claros: la apelación no es un nuevo juicio, la prueba en segunda instancia es excepcional, los recursos diferidos deben retomarse y el escrito debe estar firmado. A la vez, el cobro de los honorarios es parte esencial del ejercicio profesional: toda regulación es apelable en 5 días, fundar es una facultad con ventajas y riesgos, la cámara tiene 30 días para resolver y el pago debe hacerse en 10 días desde la firmeza para no incurrir en mora e intereses. La inteligencia artificial es una herramienta útil, pero el criterio, la estrategia y la responsabilidad profesional siguen siendo del abogado.
} \\
\bottomrule
\end{longtable}

\end{document}
